\documentclass[10pt]{article}

% —— Core packages ——
\usepackage{amsmath,amssymb}
\usepackage{stmaryrd}
\usepackage[utf8]{inputenc}
\usepackage{textcomp}
\usepackage{times}
\usepackage{microtype}
\emergencystretch=1em
\usepackage[most]{tcolorbox}
\tcbuselibrary{theorems,skins,breakable}

% —— Page layout ——
% Compact body, still preserving right margin for handwritten notes.
\usepackage[
  a4paper,
  top=14mm,
  bottom=14mm,
  left=12mm,
  right=62mm
]{geometry}

% Extra header width so the title/rule spans into the note area.
% = right margin (62mm) - left margin (12mm) = 50mm.
\newlength{\notesmargin}
\setlength{\notesmargin}{50mm}

% —— Modal logic and Greek ——
\DeclareUnicodeCharacter{03B9}{\ensuremath{\iota}}
\DeclareUnicodeCharacter{25A1}{\ensuremath{\Box}}
\DeclareUnicodeCharacter{25C7}{\ensuremath{\Diamond}}
\DeclareUnicodeCharacter{03BB}{\ensuremath{\lambda}}
\DeclareUnicodeCharacter{03C6}{\ensuremath{\varphi}}
\DeclareUnicodeCharacter{00AC}{\ensuremath{\neg}}
\DeclareUnicodeCharacter{2227}{\ensuremath{\wedge}}
\DeclareUnicodeCharacter{22A2}{\ensuremath{\vdash}}

\newenvironment{definitionbox}[1][]{%
  \begin{center}
  \begin{tcolorbox}[
    enhanced,
    breakable,
    width=0.9\textwidth,
    colback=gray!2,
    colframe=gray!60,
    coltitle=black,
    fonttitle=\bfseries\small,
    title={Definition: #1},
    boxed title style={
      colback=gray!3,
      colframe=black!60,
      sharp corners,
      boxrule=0.4pt,
    },
    boxrule=0.4pt,
    sharp corners,
    left=4pt,right=4pt,top=4pt,bottom=4pt,
    before skip=6pt,
    after skip=6pt,
  ]}%
  {\end{tcolorbox}\end{center}}

% —— Tables ——
\usepackage{array}
\usepackage{longtable}
\usepackage{booktabs}
\usepackage{calc}
\newcolumntype{C}[1]{>{\raggedright\arraybackslash}p{\dimexpr#1\linewidth-4\tabcolsep\relax}}

% —— Paragraph spacing ——
\usepackage{parskip}
\setlength{\parindent}{0pt}
\setlength{\parskip}{0.25em}

% —— Section formatting ——
\usepackage{titlesec}
\usepackage{needspace}
\usepackage{etoolbox}

\setcounter{secnumdepth}{3}

\titleformat{\section}[hang]
  {\normalfont\large\bfseries}
  {\thesection.}{0.7em}{}

\titleformat{\subsection}[hang]
  {\normalfont\normalsize\bfseries}
  {\thesubsection.}{0.7em}{}

\titleformat{\subsubsection}[hang]
  {\normalfont\small\bfseries}
  {\thesubsubsection.}{0.7em}{}

\titlespacing*{\section}{0pt}{1.0em}{0.35em}
\titlespacing*{\subsection}{0pt}{0.75em}{0.25em}
\titlespacing*{\subsubsection}{0pt}{0.55em}{0.15em}

% —— Optional page break before each section except first ——
% Disabled for compactness. Uncomment if you want each section on a new page.
% \newif\iffirstsection
% \firstsectiontrue
% \pretocmd{\section}{%
%   \iffirstsection\firstsectionfalse\else\clearpage\fi
% }{}{}

% —— Lists ——
\usepackage{enumitem}
\setlist{
  nosep,
  leftmargin=1.5em,
  topsep=0.15em,
  partopsep=0pt,
  parsep=0pt,
  itemsep=0.05em
}
\setlist[itemize,1]{label=\textbullet}
\setlist[itemize,2]{label=\textendash}

% —— Figures ——
\usepackage{graphicx}
\usepackage{float}
\setkeys{Gin}{keepaspectratio,width=0.65\linewidth}
\makeatletter
\def\fps@figure{H}
\makeatother
\graphicspath{{figures/}}

% —— URL/DOI wrapping ——
\usepackage{xurl}

% —— CSL bibliography ——
\newlength{\cslhangindent}
\setlength{\cslhangindent}{1.2em}
\newenvironment{CSLReferences}[2]%
  {\begin{enumerate}[label={[\arabic*]},leftmargin=\cslhangindent,itemsep=0pt,parsep=0pt,topsep=0pt]%
   \raggedright\sloppy
   \renewcommand{\bibitem}[2][]{\item}}%
  {\end{enumerate}}
\providecommand{\CSLLeftMargin}[1]{\item[#1]}
\providecommand{\CSLRightInline}[1]{#1}

% —— Hyperref ——
\usepackage[colorlinks=true, linkcolor=black, urlcolor=blue, citecolor=blue]{hyperref}
\urlstyle{same}

% —— Fix for Pandoc/longtable "none" counter ——
\makeatletter
\@ifundefined{c@none}{%
  \newcounter{none}%
}{}
\makeatother
\renewcommand{\thenone}{}
\providecommand*{\theHnone}{\arabic{none}}

% —— Line spacing and general flow ——
\linespread{1.08}
\raggedbottom

% —— Block quotes ——
\renewenvironment{quote}{%
  \par\vskip 0.25em
  \leftskip=1.5em\rightskip=1.5em\small\itshape
  \noindent\ignorespaces
}{\par\vskip 0.25em}

% —— Pandoc compatibility ——
\providecommand{\tightlist}{}
\providecommand{\pandocbounded}[1]{#1}

% —— Citeproc fixes ——
\makeatletter
\providecommand{\citeproctext}{}
\let\@oldprotect\protect
\def\@citeprocprotection{\let\protect\@oldprotect}
\newcommand{\citeproc}[2]{#2}
\protected\def\citeprocfootnote#1{\footnote{#1}}
\makeatother

% —— Special character definitions ——
\DeclareUnicodeCharacter{2203}{\ensuremath{\exists}}
\DeclareUnicodeCharacter{2200}{\ensuremath{\forall}}
\DeclareUnicodeCharacter{2228}{\ensuremath{\vee}}
\DeclareUnicodeCharacter{2192}{\ensuremath{\rightarrow}}
\DeclareUnicodeCharacter{2194}{\ensuremath{\leftrightarrow}}
\DeclareUnicodeCharacter{2234}{\ensuremath{\therefore}}
\DeclareUnicodeCharacter{2248}{\ensuremath{\approx}}
\DeclareUnicodeCharacter{2208}{\ensuremath{\in}}
\DeclareUnicodeCharacter{2260}{\ensuremath{\neq}}
\DeclareUnicodeCharacter{2713}{\checkmark}
\DeclareUnicodeCharacter{2212}{-}

% —— Angle brackets for LPCs ——
\newcommand{\llangle}{\langle\!\langle}
\newcommand{\rrangle}{\rangle\!\rangle}
\DeclareUnicodeCharacter{27EA}{\ensuremath{\llangle}}
\DeclareUnicodeCharacter{27EB}{\ensuremath{\rrangle}}

% —— Custom commands for frequent notation ——
\newcommand{\Rfr}[2]{\ensuremath{\mathcal{R}(#1, #2)}}
\newcommand{\LPC}[1]{\ensuremath{\llangle#1\rrangle}}

% —— Rule colour ——
\usepackage[dvipsnames]{xcolor}
\definecolor{ruleColor}{gray}{0.60}

% ============================================================
\begin{document}
% ============================================================

% —— Header block: spans body + notes margin ——
\noindent\begin{minipage}{\dimexpr\textwidth+\notesmargin\relax}

  \noindent{\sffamily\Large\bfseries Chapter 2 --- Identity, Objects, and Designators}%
  \hfill{\sffamily\normalsize 2026-04-19}

    \vspace{2pt}
  \noindent{\sffamily\small\itshape Reframed: setup and positioning for the cognitive-mechanism gap}
  
  \vspace{4pt}
  \noindent{\color{ruleColor}\rule{\linewidth}{0.35pt}}

\end{minipage}

\vspace{3pt}

% —— Body: compact left column, wide right margin for notes ——
\section{Chapter 2 --- Identity, Objects, and Designators}\label{chapter-2-identity-objects-and-designators}

Kripke's necessity-of-identity argument is widely taught in a single line: if two designators are rigid and co-refer, then necessarily they co-refer. The formula \(\forall x \forall y (x=y \to \Box(x=y))\) delivers the modal consequence; rigid designation is the substantive premise that makes the consequence reach ordinary names. This chapter takes that reading as its starting point rather than its conclusion. The chapter's task is not to argue that rigidity is presupposed --- that is consensus --- but to specify what kind of presupposition it is, what the formula does and does not do once rigidity is granted, and what would be required to discharge the presupposition rather than merely posit it. The last question is what the rest of the dissertation answers.

\subsection{The Formula and the Standard Reading}\label{the-formula-and-the-standard-reading}

The necessity-of-identity formula was first derived by Ruth Barcan Marcus in 1947 and presented in simplified form by Quine in 1953. Kripke's ``Identity and Necessity'' (1971) and \emph{Naming and Necessity} (1980) gave it the philosophical framing under which it is now standardly read. In \emph{Naming and Necessity}, Kripke is explicit that the principle, applied to ordinary identity claims, derives directly once rigid designation is assumed: if `a' and `b' are rigid designators, and `a = b' is true, then `a = b' is necessarily true.\footnote{Kripke, \emph{Naming and Necessity} (1980), pp.~100--105. The conditional structure is also visible in Kripke's earlier ``Identity and Necessity'' (1971), where the formula is presented and then immediately illustrated with the Benjamin Franklin case, with the rigidity of ``Benjamin Franklin'' doing the substantive work.} The argument has two parts: a formal core (the modal logic of identity under fixed assignment) and a metasemantic premise (that ordinary proper names function as rigid designators). The formal core delivers the modal consequence cleanly; the metasemantic premise carries the philosophical weight.

The standard reading separates these two parts. Speaks's lecture notes present the argument as a numbered derivation in which rigidity appears as an explicit premise; if either of the two terms fails to be a rigid designator, the derivation does not go through.\footnote{Speaks, ``Kripke on the necessary a posteriori I: identity sentences'' (lecture notes, McGill PHIL 415), §2: ``Take any identity sentence \(n=m\), where \(n\) and \(m\) are both rigid designators. Suppose that the sentence is true. It then seems to follow that it is also necessarily true\ldots{}''} The Stanford Encyclopedia entry on rigid designators treats the rigidity of names as Kripke's substantive thesis and the modal consequence as following from it together with the assumption that the identity is true.\footnote{LaPorte, ``Rigid Designators,'' \emph{Stanford Encyclopedia of Philosophy}, §1.1.} The Wikipedia entry on the necessity of identity, summarizing the consensus, presents Kripke's later strategy as deriving the principle ``directly, assuming what he called rigid designation.''\footnote{``Necessity of identity,'' Wikipedia (consulted 2026): ``In the later \emph{Naming and Necessity}, Kripke suggested that the principle could be derived directly, assuming what he called rigid designation.''} Burgess's careful technical reconstruction of the derivation in \emph{Synthese} makes the conditional structure explicit, distinguishing the source, status, and significance of the formal result from the substantive metasemantic claims it is used to support.\footnote{Burgess (2014), ``On a Derivation of the Necessity of Identity,'' \emph{Synthese} 191, 1567--85.}

So the chapter's headline observation --- that the formula does not by itself establish rigidity for names --- is not a discovery. It is the standard reading, and Kripke himself frames the argument in these terms. The substantive case for the rigidity of names is made elsewhere, on intuitive grounds: the counterfactual cases of \emph{Naming and Necessity} (Nixon could have lost the 1968 election; Aristotle could have died in infancy and never taught Alexander), and the Gödel/Schmidt thought experiment (the name ``Gödel'' would still refer to Gödel even if everything we attribute to him had in fact been done by someone else).\footnote{Kripke, \emph{Naming and Necessity}, pp.~74--78 (Nixon), pp.~60--62 and 75--76 (Aristotle), pp.~83--87 (Gödel/Schmidt). The cases function together as the modal-intuition battery against descriptivism.} These are where the rigid-designation thesis lives. The formula provides the modal closure once the thesis is granted.

What this chapter does is sit at the joint between the formal core and the metasemantic premise, and ask a more specific question: what does the metasemantic premise commit us to, and what would discharging it cognitively require? The dissertation's contribution begins where the standard reading leaves off --- not in disputing the conditional, but in supplying the mechanism by which the antecedent of the conditional gets satisfied for actual speakers using actual names.

\subsection{What the Formula Actually Does}\label{what-the-formula-actually-does}

The formal core of the argument is uncontroversial. Under any assignment function \(g\) in a Kripke model, if \(g(x) = g(y)\), then in every world the model constructs, \(g(x) = g(y)\). The assignment function does not vary across worlds in the standard semantics; once an assignment is fixed, the variables continue to pick out their assigned values across every world the operator \(\Box\) quantifies over. The formula is valid by construction.\footnote{Schwarz (2006), ``The Necessity of Identity: A Circularity Objection,'' identifies a concern with the Leibniz's Law step in some presentations of the derivation. The concern is downstream of the present analysis: the formal validity at the level of variables under assignment is not in question; what is in question is what happens when one tries to extend the result to intensional cases. Garson (2024), ``Logics for Rigidity,'' in Weiss \& Birman (eds.), \emph{Saul Kripke on Modal Logic} (Springer), pp.~193--208, surveys the technical situation and finds that the QML systems that validate the necessity of identity for rigid designators are ``very different from ones Kripke developed,'' with ``serious costs whichever solution we attempt.''}

Three observations follow. They are unremarkable in themselves; together they specify what work the formula can and cannot do for the philosophical project.

The formula's validity rests on the rigidity of variables under assignment, which is a feature of the formal apparatus, not a feature of natural-language names. When Kripke applies the formula to names, he is applying a formal result whose validity holds for one kind of designator (variables under assignment) to another kind of designator (proper names). The application requires the assumption that names exhibit the same world-invariance variables have by construction. This is the rigid-designation thesis itself, not a consequence of it.

The formula has no resources for the intensional dimension of the cases that motivate Kripke's project. Hesperus/Phosphorus, Benjamin Franklin and the inventor of bifocals, water and H\(_2\)O --- these are cases where identity is empirically discovered and then recognized as necessary. The formula regiments the modal status once rigidity is granted; the epistemic route by which speakers reach the identity, and the cognitive capacities involved in using a name as rigid in the first place, lie outside what the formula can address.\footnote{A constraint from the \emph{Tractatus} sharpens what kind of identity claim is at stake. Wittgenstein observes (5.5303) that ``to say of two things that they are identical is nonsense, and to say of one thing that it is identical with itself is to say nothing at all.'' Identity claims are claims about designators co-referring, not about objects themselves being one. Kripke's formula, applied to names, is consistent with this: it concerns the modal status of co-reference between designators, not a metaphysical relation among objects. The Tractarian connection is surprisingly underdeveloped in Kripke commentary, which has concentrated on rule-following (Kripke 1982) and on the technical literature on identity-elimination (Wehmeier 2008, 2012; Lampert and Säbel 2016) rather than on the \emph{Tractatus}'s bearing on the necessity-of-identity argument.}

The formula does not adjudicate between competing accounts of how names achieve the world-invariance it presupposes. Kripke's own account, developed in \emph{Naming and Necessity}, runs through baptism and causal-historical transmission. Other accounts --- descriptivist, hybrid causal-descriptivist, Mental Files --- give different stories about what makes names function as rigid designators. The formula is silent on the choice. It tells us what follows if names are rigid; the substantive metasemantic question is what makes them so. Kripke himself distinguishes the rigidity of names --- \emph{de jure} rigidity, built into the kind of designator a name is --- from \emph{de facto} rigidity, which attaches to descriptions whose satisfiers happen to be necessarily fixed (such as ``the least prime'').\footnote{Kripke, \emph{Naming and Necessity}, p.~21, fn. 21. The \emph{de jure} / \emph{de facto} distinction matters here because the formula's presupposition is \emph{de jure} rigidity --- rigidity by how the designator functions, not by metaphysical accident.} He is also explicit that the wide-scope behavior of descriptions in modal contexts is not what the rigid-designation thesis is asserting: the thesis is about the kind of designator names are, not about scope manipulation.\footnote{Kripke, \emph{Naming and Necessity}, p.~12, fn. 15. Kripke notes that taking a description outside the scope of a modal operator can produce something that behaves like rigidity, and is explicit that this is not what the rigid-designation thesis is asserting. The point is later developed by Kaplan's \emph{dthat} operator, examined in Chapter 3.}

\subsection{Three Levels of ``Object''}\label{three-levels-of-object}

The standard reading distinguishes the formal core from the metasemantic premise. The chapter adds a finer distinction within the metasemantic premise itself. When the formula is applied to ordinary identity claims, the word ``object'' --- used in Kripke's own gloss and throughout the secondary literature --- covers three distinct things, and the philosophical work happens at a different level than the formal validity does.

\emph{Domain elements} are mathematical objects: points in a set, self-identical by set-theoretic construction, their identity invariant across any world the model constructs. The formula's validity operates at this level. Variables are assigned to domain elements; the assignment function is the device through which formal rigidity is achieved.

\emph{Encountered particulars} are individuals as they figure in perception and re-identification: this dog here, that star there, the man one might have met. They are individuated through encounter; their identity across encounters is an epistemic achievement, with room for error. The early astronomer saw Hesperus in the evening sky and Phosphorus in the morning sky; the identity of the two had to be established, not given with the encounters. Encountered particulars are the bridge between the formal apparatus and the philosophical applications, because they have the concreteness that descriptions can apply to (so the cases Kripke discusses are intelligible at all) while also being the sort of thing that can bear a name.

\emph{Named individuals} are bearers of proper names: Aristotle, Venus, Hesperus, Benjamin Franklin. They are individuated through baptism and tracked across absence through testimonial networks. Their identity across different names is discoverable --- the Hesperus/Phosphorus case is the paradigm. The philosophical thesis of rigid designation operates at this level: names hold their referents fixed across counterfactual variation, where descriptions could track different individuals in different worlds.

The Benjamin Franklin passage, which Kripke uses to illustrate the formula immediately after presenting it, exhibits the level-shift in compressed form. The passage begins with the formal apparatus (``there is an object x such that x invented bifocals''), passes through descriptive identification (``the inventor of bifocals,'' ``the first Postmaster General''), and arrives at a name (``x and y are both Benjamin Franklin''). Read at the level of variables, the descriptions do not apply; mathematical objects do not invent bifocals. Read at the level of encountered particulars, the descriptions apply but the identity claim becomes awkward in the way Wittgenstein's constraint anticipates --- identity between two encountered particulars is either nonsense or self-identity. Read at the level of names, the necessity follows from the rigidity of ``Benjamin Franklin'' --- but only because rigidity has been imported, not because the formula has delivered it.

What the passage does is move smoothly between these levels in English. The smoothness is rhetorical, not logical. Each level handles part of the argument; no single level handles the whole. The metasemantic premise --- that ordinary names exhibit the world-invariance variables have by construction --- is what makes the move from the first level to the third coherent. The premise is supplied by the rigid-designation thesis, which Kripke argues for separately on intuitive grounds (the modal-intuition tests of \emph{Naming and Necessity}, the counterfactual cases about Nixon and Aristotle).\footnote{One could trace the four readings of ``object x'' --- mathematical, concrete, positional, name-substitute --- and show that no single reading captures the whole sequence of the passage. The exercise yields a typology consistent with the level-shift analysis above, but the typology adds little to the underlying observation that the passage moves between levels without committing to any one. The reader interested in the more detailed reading-by-reading analysis is referred to the working notes.}

The three-level distinction is useful for the rest of the dissertation. The cognitive question --- what enables speakers to use names as rigid designators --- is a question about how the metasemantic premise gets satisfied at the level of named individuals. Local rigidity, developed in Chapters 4 and 5, is the proposed mechanism. It operates between encountered particulars and named individuals: it explains how reference to an encountered particular can be lifted, through baptism, into a name that tracks the referent across absence. The formal apparatus does not need a cognitive story; the philosophical application does.

\subsection{The Cognitive-Mechanism Gap}\label{the-cognitive-mechanism-gap}

What the standard reading leaves open is the cognitive question. Kripke argues that names are rigid; he does not argue that they are rigid because of any particular cognitive mechanism. The causal-historical sketch in \emph{Naming and Necessity} --- baptism fixes the reference, the name propagates through chains of communication --- is offered explicitly as a ``better picture'' rather than as a theory.\footnote{The ``better picture'' framing is widely cited in the secondary literature; see, e.g., S. Marc Cohen's lecture notes (University of Washington) and Funkhouser (PHIL 4233 lecture notes), both of which foreground Kripke's self-description as a picture rather than a theory.} Kripke is unusually direct about this self-positioning: he is not, he says, ``trying to give any sort of necessary and sufficient conditions for reference\ldots{} in any rigorous sense''; he is ``painting a general picture.''\footnote{Kripke, \emph{Naming and Necessity}, pp.~93--94. The full passage: ``I'm not, in fact, trying to give any sort of necessary and sufficient conditions for reference\ldots{} in any rigorous sense'' (p.~94); ``I think it is a much better picture than the picture presented by a description'' (p.~93).} The picture describes the output of the relation between names and their referents; the cognitive mechanism that produces and sustains the output is not analyzed.

That a cognitive mechanism is needed at all, and that the available candidates may be too narrow, is something Kripke himself notices in passing. In ``Identity and Necessity,'' he quotes Russell at length on demonstratives as the only true names --- terms that ``really just name an object without describing it'' --- and on acquaintance as the only ground secure enough for identity claims to be ``exempt from all imaginable doubt.'' Kripke's gloss is that on Russell's view, only names of immediate sense data, expressed by demonstratives like ``this'' and ``that,'' qualify as genuine names.\footnote{Kripke, ``Identity and Necessity'' (1971), p.~144. The Russell passage Kripke quotes is from \emph{The Philosophy of Logical Atomism} (1918) and runs: ``if we want to reserve the term `name' for things which really just name an object without describing it, the only real proper names we can have are names of our own immediate sense data, objects of our own immediate acquaintance. The only such names which occur in language are demonstratives like `this' and `that.'\,''} The passage is not used by Kripke to endorse Russell's account; it is used to motivate the contrast with ordinary proper names, which are rigid in Kripke's sense but do not meet Russell's Cartesian standard. What the passage acknowledges, without developing, is that there is a candidate cognitive mechanism for rigidity --- demonstrative grounding under acquaintance --- and that Russell himself made it too narrow to cover the cases Kripke wants to cover. The mechanism Russell identified is the right kind of thing; the scope at which Russell allowed it is too restricted. Closing this gap --- preserving Russell's mechanism while extending its scope --- is the dialectical move that Chapters 4 and 5 will make explicit.

This is acknowledged in the secondary literature. The Stanford Encyclopedia's discussion of rigid designation notes that Kripke's appeal to stipulation in the treatment of \emph{de jure} rigidity is ``more like a promissory note than the satisfaction of an explanatory obligation.''\footnote{LaPorte, ``Rigid Designators,'' §3.} The promissory note is for an account of the mechanism by which names achieve, in actual linguistic practice, what variables achieve by construction. The literature has attempted to discharge the promissory note in several ways. Evans (1982) defends a ``discriminating conception'' requirement on reference. Recanati (2012) develops the account into the Mental Files architecture. Devitt (1981) offers a causal-descriptivist hybrid that supplements the causal chain with a descriptive component. Each addresses part of what is needed; Chapter 3 examines them in detail.

The dissertation's contribution begins here. Chapter 3 surveys the existing cognitive accounts and argues that each identifies the right desideratum but stops short of supplying the mechanism. Chapter 4 introduces the analytical vocabulary --- the \(x^o\)/\(x^c\) mode distinction --- that makes the missing mechanism visible. Chapter 5 develops the positive account: local rigidity as the cognitive primitive, baptism as the lifting of local rigidity into a socially portable name, the causal chain as the testimonial network that sustains the lifting across absence. The metasemantic premise that the standard reading correctly identifies as substantive --- that ordinary names function as rigid designators --- gets, on the proposed account, a cognitive grounding it has not had.

What this chapter has done, then, is set the stage. The headline observation is consensus, sourced to Kripke and the standard literature. The three-level distinction adds a finer specification of where the cognitive question arises. The cognitive-mechanism gap is identified as the place where the dissertation makes its substantive contribution. The chapters that follow take up the gap and propose a way to close it.

\vfill
\noindent\rule{\textwidth}{0.35pt}\\[1pt]
{\tiny florin.cojocariu@s.unibuc.ro, \today}

\end{document}